\chapter*{译者说明}
\addcontentsline{toc}{chapter}{译者说明}
\markboth{译者说明}{译者说明}

\section*{底本}

本书译自 A.~Beilinson 与 V.~Drinfeld 的专著 \emph{Quantization of Hitchin's Integrable
System and Hecke Eigensheaves}。该书没有正式出版本，流传的是作者自 1990 年代起陆续修订的
pdf\TeX{} 稿；本译本依据的是 2020 年 11 月的重排版，共 386 页，无目录、无书签。

原书页码与 PDF 页序一一对应（第 1 页即标题页），因此译本中每页对应的原文位置可以直接用页码
指认。各章在原书中的起始页为：第 0 章 p.3，第 1 章 p.6，第 2 章 p.15，第 3 章 p.61，
第 4 章 p.122，第 5 章 p.192，第 6 章 p.230，第 7 章 p.236，第 8 章 p.367，第 9 章 p.370，
第 10 章 p.375，参考文献 p.381。

\section*{编号}

原书的编号制度有三层，译本一律照抄，不重新编号：

\begin{itemize}
  \item 章、节沿用 \texttt{章.节} 形式，例如 2.7 表示第 2 章第 7 节；
  \item 编号段落用 \texttt{章.节.序号} 形式，例如 2.6.7；这类编号直接取自原书，
        不依赖 \LaTeX{} 计数器，故增删段落不会破坏编号；
  \item 公式按章连续编号，例如 (26) 表示第 2 章第 26 式。
\end{itemize}

正文中引用原书编号时写成「式（26）」「2.6.3」「见 3.1.13」，不使用交叉引用命令，
以保证与原文一致。原书自身的坏引用（印成 \texttt{???} 或指向错误编号处）照录，并在该处
用批注说明实际所指。

\section*{页码锚点}

每一页原文在译本中对应一个不可见的页码锚点，宏名为 \verb|\SourcePage{N}|，其中 $N$ 是
原书页码。校对时用 \texttt{orig:N} 作超链接目标即可跳到该页开头。

\section*{批注}

译本中的批注一律写成 \LaTeX{} 注释行，形如

\begin{verbatim}
% [Tutor Note: 原书 p.65 此处印 $\omega_X^{\otimes 2}$，
%  但按 3.1.2 的权计算应为 $\omega_X^{\otimes 3}$；译文照录原书，未擅改。]
\end{verbatim}

批注只写三类内容：易错预警（读者容易看错或算错之处）、含糊点拨（原文表述含混之处）、
背景直觉（补一句为什么这样定义）。原书的笔误一律照录并加批注，译本不作数学上的擅自修正；
确需给出更正确读法时，写在批注里，并标明依据。

\section*{术语}

术语以工程文件中的术语表为准，开译前冻结，逐章统一。少数词采取不译或首现括注的处理：
\emph{oper} 一词在几何朗兰兹文献中已成专名，全书不译；\emph{torsor} 首现时括注「主齐性空间」，
其后不译；\emph{Pfaffian} 首现时括注「普法夫式」。

\section*{公式}

原书 PDF 的文字层把数学压平（上下标丢失、求和号与映射箭头错乱），因此译本中每一处公式、
定义与定理陈述都对照页面图像逐符号核对后重排，不依赖文字层。

\section*{致谢}

底本由作者提供并授权流传。译本的排版沿用英文原稿的编号与页面结构，便于与原文对照阅读。
